% Set up the document class and basic geometry
\documentclass[12pt, a4paper]{article}
\usepackage[utf8]{inputenc}
\usepackage[T1]{fontenc}
\usepackage[polski]{babel}
\usepackage{geometry}
\geometry{margin=2.5cm}
\usepackage{hyperref}
\usepackage{enumitem}
\usepackage{titlesec}

% Define the title section
\title{\textbf{Dokument Początkowy Projektu}\\ \large Przedmiot: Algorytmy Wspomagania Decyzji}
\author{Mariusz Gąsiewicz, Filip Ferenc}
\date{\today}

\begin{document}

\maketitle

\section{Temat projektu}
Inteligentny algorytm adaptacji łącza (Link Adaptation / MCS Selection) dla stacji bazowych 5G (gNB) z wykorzystaniem modeli uczenia maszynowego i destylacji wiedzy.

\section{Cel projektu}
Głównym celem projektu jest opracowanie algorytmu wspomagania decyzji, który w czasie rzeczywistym dobiera optymalny schemat modulacji i kodowania (MCS -- \textit{Modulation and Coding Scheme}) dla użytkowników w sieci 5G. 

W przeciwieństwie do tradycyjnych metod (np. OLLA -- \textit{Outer Loop Link Adaptation}) opartych na prostych heurystykach, proponowane rozwiązanie ma skuteczniej radzić sobie w dynamicznych środowiskach radiowych. Algorytm będzie dążył do maksymalizacji przepustowości (Throughput) przy jednoczesnym rygorystycznym utrzymaniu wskaźnika błędów paczek (BLER -- \textit{Block Error Rate}) poniżej zadanego progu 10\%, niezależnie od warunków kanałowych, prędkości użytkownika czy poziomu interferencji (SIR).

\section{Wykorzystywane algorytmy i podejście}
Projekt opiera się na zaawansowanym przetwarzaniu i modelowaniu decyzji z wykorzystaniem następujących koncepcji:

\begin{itemize}
    \item \textbf{Ordinal Regression (Regresja porządkowa):} Indeksy MCS mają naturalny porządek -- wyższy MCS oznacza większą przepustowość, ale mniejszą odporność na błędy. Problem decyzyjny zostanie zdekomponowany na zestaw binarnych klasyfikatorów z użyciem algorytmu \textit{HistGradientBoosting}. Zastosowane zostaną monotoniczne więzy obostrzeń (\textit{monotonic constraints}), aby zagwarantować fizyczny sens podejmowanych decyzji.
    \item \textbf{Destylacja wiedzy (Knowledge Distillation):} Z uwagi na krytyczne wymagania opóźnieniowe w warstwie MAC (rzędu mikrosekund), złożony model złożeń nie może działać w czasie rzeczywistym. Decyzje modelu nadrzędnego zostaną wydestylowane do surogatowej struktury Drzewa Decyzyjnego (\textit{Decision Tree}) o ograniczonej głębokości (np. \texttt{max\_depth=10}).
    \item \textbf{Asymetryczna funkcja kosztu (Asymmetric Rank Cost):} Algorytm uwzględni asymetrię ryzyka w telekomunikacji. Przeszacowanie możliwości łącza skutkuje retransmisją HARQ (wysoki koszt), natomiast niedoszacowanie prowadzi jedynie do suboptymalnej przepustowości.
    \item \textbf{Safety Fallback:} System decyzyjny ML zostanie zabezpieczony klasycznym algorytmem OLLA (sterowanie marginesem) oraz mechanizmem awaryjnym (kill-switch) reagującym na serię błędów dekodowania.
\end{itemize}

\section{Narzędzia i technologie}
\begin{itemize}
    \item \textbf{NVIDIA Sionna (TensorFlow):} Zaawansowany symulator warstwy fizycznej (PHY) 5G. Posłuży do generacji realistycznych zbiorów danych uwzględniających m.in. siatkę OFDM, kodowanie LDPC, modele kanałów 3GPP i efekt Dopplera.
    \item \textbf{Python (Scikit-learn, Pandas, NumPy):} Główne środowisko do analizy danych, inżynierii cech (np. wprowadzanie szumu estymacji kanału) oraz trenowania modeli uczenia maszynowego.
    \item \textbf{C++:} Język docelowy do wdrożenia wyuczonej polityki decyzyjnej (drzewa), zoptymalizowanej pod kątem integracji z otwartymi schedulerami sieciowymi (np. srsRAN, FlexRIC).
\end{itemize}

\section{Uzasadnienie naukowe i biznesowe}
Wybór tematu jest ściśle powiązany z profilem przedmiotu. Projekt modeluje rygorystyczne środowisko decyzyjne, w którym kluczowa jest asymetryczna funkcja kosztu (kary za opóźnienia), a nie tylko surowa dokładność (Accuracy). Uwzględnia również krytyczny czas podjęcia decyzji (\textit{inference latency}) poprzez zastosowanie destylacji modeli, a wykorzystanie obostrzeń monotonicznych stanowi elegancki sposób integracji wiedzy domenowej ze strukturą decyzyjną maszyny.

\section{Harmonogram projektu (Roadmapa)}
Poniższy plan określa główne etapy realizacji projektu:

\begin{itemize}
    % Phase 1: Data preparation
    \item \textbf{Etap 1: Konfiguracja symulatora i generacja danych}
    \begin{itemize}
        \item Konfiguracja środowiska NVIDIA Sionna.
        \item Wygenerowanie zestawów danych treningowych i testowych z uwzględnieniem różnych modeli kanałów i zjawisk fizycznych.
    \end{itemize}
    
    % Phase 2: Feature engineering and base model
    \item \textbf{Etap 2: Inżynieria cech i trenowanie modelu nadrzędnego}
    \begin{itemize}
        \item Przygotowanie potoku przetwarzania danych w Pythonie.
        \item Wdrożenie modelu \textit{HistGradientBoosting} z uwzględnieniem założeń regresji porządkowej i więzów monotonicznych.
    \end{itemize}
    \newpage
    % Phase 3: Distillation and cost function
    \item \textbf{Etap 3: Zastosowanie asymetrycznej funkcji kosztu i destylacja}
    \begin{itemize}
        \item Optymalizacja progów decyzyjnych modelu pod kątem rygorystycznego wymogu BLER $< 10\%$.
        \item Przeprowadzenie procesu destylacji wiedzy do lekkiego modelu opartego na Drzewie Decyzyjnym.
    \end{itemize}
    
    % Phase 4: Safety constraints and C++ translation
    \item \textbf{Etap 4: Warstwy bezpieczeństwa i eksport logiki}
    \begin{itemize}
        \item Zaprojektowanie i symulacja mechanizmu \textit{Safety Fallback} (OLLA, kill-switch).
        \item Konwersja reguł Drzewa Decyzyjnego na optymalny kod C++ (\texttt{if/else}).
    \end{itemize}
    
    % Phase 5: Final evaluation
    \item \textbf{Etap 5: Testy i dokumentacja końcowa}
    \begin{itemize}
        \item Przeprowadzenie testów wydajnościowych rozwiązania.
        \item Porównanie efektywności modelu z klasycznymi heurystykami.
        \item Opracowanie raportu końcowego.
    \end{itemize}
\end{itemize}

\end{document}